\section{METHODOLOGY}
\label{sec:methodology}
% =============================================================================

\subsection{Dataset Overview}

Tập dữ liệu sử dụng trong nghiên cứu này là tập dữ liệu Tiêu thụ Điện năng Ngành thép~\cite{uci2023}, có sẵn công khai từ Kho lưu trữ Học máy UCI.
Nó chứa 35.040 quan sát được ghi lại tại các khoảng 15 phút từ một nhà máy sản xuất thép Hàn Quốc trong năm 2018.
Mỗi quan sát bao gồm 11 biến ban đầu: công suất tác dụng ($P$), công suất phản kháng trễ và sớm ($Q_{\text{lag}}$, $Q_{\text{lead}}$), hệ số công suất (PF), lượng khí CO$_2$ tương đương, loại tải, trạng thái tuần, và các biến thời gian (NSM, ngày trong tuần).

\subsection{Data Preprocessing}

Giai đoạn tiền xử lý bao gồm kiểm tra dữ liệu, loại bỏ trùng lặp, và xử lý giá trị thiếu (nếu có) bằng nội suy tuyến tính.
Đối với dữ liệu chuỗi thời gian, nội suy tuyến tính được áp dụng:
\begin{equation}
y(t) = y(t_1) + \frac{y(t_2) - y(t_1)}{t_2 - t_1}(t - t_1)
\label{eq:interp}
\end{equation}

Một vấn đề quan trọng được phát hiện trong quá trình audit dữ liệu là cột Power Factor gốc được lưu dưới dạng phần trăm (0--100) thay vì hệ số (0--1).
Việc không phát hiện lỗi này sẽ làm sai lệch hoàn toàn các tính toán vật lý phụ thuộc PF, bao gồm góc lệch pha $\varphi = \arccos(\text{PF})$.
Pipeline tự động phát hiện và scale về $[0,1]$ bằng cách kiểm tra $\max(\text{PF}) > 1{,}0$.

\subsubsection{Missingness Mechanism Analysis}

Trong quá trình thiết kế pipeline, chúng tôi tích hợp module phân tích cơ chế thiếu dữ liệu theo khung lý thuyết của Rubin~\cite{rubin1976} như một cơ chế audit phòng ngừa.
Cơ chế thiếu dữ liệu được phân loại thành ba nhóm: MCAR (Missing Completely At Random), MAR (Missing At Random), và MNAR (Missing Not At Random).
Việc xác định đúng cơ chế ảnh hưởng trực tiếp đến lựa chọn phương pháp imputation và độ tin cậy của suy luận thống kê sau này.

\textbf{Little's MCAR Test (Đa biến xấp xỉ).}
Với mỗi cột có missing, chúng tôi chia mẫu thành nhóm missing và nhóm observed, sau đó so sánh vector trung bình của các cột số khác bằng kiểm định $t$-test độc lập với hiệu chỉnh Bonferroni.
Nếu không có sự khác biệt ý nghĩa nào, ta không thể bác bỏ giả thuyết MCAR:
\begin{equation}
H_0: \boldsymbol{\mu}_{\text{obs}} = \boldsymbol{\mu}_{\text{mis}} \quad \text{(MCAR)}
\label{eq:mcar}
\end{equation}

\textbf{MAR Logistic Test.}
Missing indicator $M_j \in \{0,1\}$ của cột $j$ được mô hình hóa bằng hồi quy logistic dựa trên các giá trị observed của các cột khác:
\begin{equation}
\log\frac{P(M_j=1)}{1-P(M_j=1)} = \beta_0 + \sum_{k \neq j} \beta_k X_k
\label{eq:mar}
\end{equation}
Nếu AUC $> 0{,}60$, missing có khả năng phụ thuộc vào dữ liệu observed → MAR.

\textbf{MNAR Range Analysis.}
Nếu tỷ lệ missing tập trung bất thường ở các bin giá trị thấp hoặc cao của chính cột đó, dấu hiệu MNAR xuất hiện.
Điều này thường gặp trong hệ thống cảm biến khi thiết bị ngừng hoạt động ở vùng giá trị cực đoan.

\subsubsection{Data Assertions}

Để đảm bảo tính tái lập và phát hiện lỗi tự động, một module kiểm tra khẳng định (data assertions) được tích hợp sau mỗi bước pipeline.
Các khẳng định này kiểm tra ràng buộc vật lý, tính nhất quán thời gian, và tính toàn vẹn dữ liệu.

Các kiểm tra chính bao gồm:
\begin{itemize}
    \item \textbf{Ràng buộc vật lý:} $P \ge 0$, PF $\in [0,1]$, $S^2 \approx P^2 + Q^2$.
    \item \textbf{Nhất quán thời gian:} Dữ liệu sắp xếp tăng dần, không trùng lặp timestamp, NSM khớp với \texttt{date}.
    \item \textbf{Toàn vẹn:} Không có giá trị null sau khi nội suy.
    \item \textbf{Tỷ lệ bất thường:} Tổng tỷ lệ anomaly $< 10\%$.
    \item \textbf{Số lượng đặc trưng:} 69 cột trong gold dataset sau feature engineering và labeling.
\end{itemize}

\begin{table}[htbp]
\caption{CÁC KHẲNG ĐỊNH DỮ LIỆU THEO GIAI ĐOẠN PIPELINE}
\label{tab:assertions}
\centering
\small
\begin{tabular}{@{}p{2.0cm}p{5.0cm}@{}}
\toprule
Giai đoạn & Các khẳng định được kích hoạt \\
\midrule
Sau làm sạch & $P \ge 0$, PF $\in [0,1]$, thời gian tăng dần, không null \\
Sau đặc trưng thời gian & Kiểm tra lại như trên \\
Sau đặc trưng vật lý & Thêm $S^2 \approx P^2 + Q^2$ \\
Sau gán nhãn & Thêm tỷ lệ anomaly $< 10\%$, schema cuối gồm 69 cột \\
Sau GAN & Thêm bảo toàn ma trận tương quan real vs synthetic \\
\bottomrule
\end{tabular}
\end{table}

\subsection{Time-Domain Feature Engineering}

Đặc trưng trễ được xây dựng với các độ trễ $k \in \{1, 2, 4, 96\}$
tương ứng với 15 phút, 30 phút, 1 giờ, và 24 giờ:
\begin{equation}
x_{\text{lag}}(k) = x(t - k)
\label{eq:lag}
\end{equation}

Thống kê cửa sổ trượt bao gồm trung bình, độ lệch chuẩn,
và độ skewness trên các cửa sổ $w \in \{24, 48, 96\}$.
Đặc biệt, rolling standard deviation tăng vọt thường là dấu hiệu báo trước sự cố quá tải hoặc chuyển đổi chế độ vận hành đột ngột.

Mã hóa chu kỳ của Số giây tính từ nửa đêm (NSM)
đảm bảo tính liên tục tại nửa đêm:
\begin{align}
\text{NSM}_{\sin} &= \sin\left(\frac{2\pi \cdot \text{NSM}}{86400}\right) \\
\text{NSM}_{\cos} &= \cos\left(\frac{2\pi \cdot \text{NSM}}{86400}\right)
\end{align}

\subsection{Frequency-Domain Feature Extraction (DWT)}

Biến đổi Wavelet rời rạc phân rã tín hiệu thành
các hệ số xấp xỉ $cA$ và chi tiết $cD$:
\begin{equation}
S(t) = cA_j(t) + \sum_{i=1}^{j} cD_i(t)
\label{eq:dwt}
\end{equation}

Wavelet mẹ $\psi(t)$ được co giãn và tịnh tiến để tạo họ hàm cơ sở:
\begin{equation}
\psi_{a,b}(t) = \frac{1}{\sqrt{a}}\psi\left(\frac{t-b}{a}\right)
\label{eq:wavelet_basis}
\end{equation}
trong đó $a$ là tham số tỷ lệ (scale) và $b$ là tham số dịch chuyển (shift).

Chúng tôi sử dụng wavelet Daubechies-4 (db4) với 3 mức phân rã.
Lựa chọn db4 dựa trên sự cân bằng giữa độ mượt (smoothness) và độ dài hỗ trợ compact:
db4 đủ mượt để tránh artifact số học, đủ ngắn để bắt giữ các sự kiện đột ngột trong tải công nghiệp.
Theo Liu et al.~\cite{liu2022modwt}, việc sử dụng Maximal Overlap DWT (MODWT)
thay vì DWT tiêu chuẩn có thể giảm thiểu rò rỉ thông tin tần số
và cải thiện độ ổn định của đặc trưng wavelet trên cửa sổ trượt.
Một cửa sổ trượt 64 mẫu (16 giờ) được sử dụng để tính toán
các đặc trưng wavelet (mean, std, energy, max\_abs) tại mỗi bước thời gian.

\subsection{Physics-Informed Feature Engineering}

Để làm phong phú tập dữ liệu với các đại lượng điện năng mang tính vật lý,
hai biến phái sinh quan trọng --- Công suất biểu kiến ($S$) và Góc lệch pha
($\varphi$) --- được tính toán từ các biến đo công suất thô.
Những đặc trưng này giúp phơi bày các chế độ vận hành bất thường mà công suất tác dụng ($P$) đơn thuần không thể phát hiện,
chẳng hạn như hiện tượng kẹt rô-to hoặc rò rỉ cách điện tiến triển,
từ đó cung cấp các chỉ báo hữu ích cho mô hình dự báo tải và phân loại rủi ro.

\textbf{1) Công suất biểu kiến ($S$):}
Công suất biểu kiến là độ lớn của công suất phức trong mạch điện xoay chiều.
Đại lượng này quyết định giới hạn chịu tải nhiệt và cơ học của dây cáp, máy biến áp và thiết bị đóng cắt.
Khi công suất tác dụng không đổi nhưng công suất phản kháng tăng (do kẹt rô-to hoặc bão hòa từ trường),
$S$ sẽ phình to --- đây là tín hiệu cảnh báo quá dòng.

\textbf{Định nghĩa:}
\begin{equation}
S = \sqrt{P^2 + Q^2}
\label{eq:apparent}
\end{equation}
Trong đó $P = \text{Usage\_kWh}$ là công suất tác dụng (kW),
và $Q = \text{Lagging\_Current\_Reactive.Power\_kVarh}$ là công suất phản kháng trễ (kVar).

\textbf{2) Góc lệch pha ($\varphi$):}
Góc lệch pha định lượng độ dịch thời gian giữa sóng điện áp và dòng điện.
Đại lượng này cung cấp thang đo tuyến tính sắc nét hơn Hệ số công suất (PF),
đặc biệt khi phụ tải trở nên phi tuyến hoặc mang tính cảm kháng cao.

\textbf{Định nghĩa:}
\begin{equation}
\varphi = \arccos(\text{PF})
\label{eq:phase}
\end{equation}
Trong đó $\text{PF} = \text{Lagging\_Current\_Power\_Factor}$ được giới hạn trong khoảng $[0, 1]$
trước khi áp dụng hàm $\arccos$ nhằm đảm bảo tính ổn định số học.
Kết quả $\varphi$ được biểu diễn bằng radian.

\begin{table*}[htbp]
\caption{CÁC ĐẶC TRƯNG MIỀN VẬT LÝ}
\label{tab:physical}
\centering
\small
\begin{tabular}{@{}p{3.0cm}p{1.2cm}p{3.0cm}p{1.2cm}p{5.4cm}@{}}
\toprule
Đặc trưng & Ký hiệu & Công thức & Đơn vị & Ý nghĩa vật lý \\
\midrule
Công suất biểu kiến & $S$ & $\sqrt{P^2+Q^2}$ & kVA & Tổng công suất đặt lên hệ thống truyền tải; cảnh báo quá tải sớm. \\
Góc lệch pha & $\varphi$ & $\arccos(\mathrm{PF})$ & rad & Thang đo tuyến tính cho tính phi tuyến của phụ tải; phân biệt rõ hơn PF gần 1. \\
\bottomrule
\end{tabular}
\end{table*}

\textbf{Giải thích vật lý và ứng dụng:}
Trong bối cảnh nhà máy thép, cả $S$ và $\varphi$ đóng vai trò chỉ báo nhạy với bất thường.
Sự gia tăng đồng thờ của $S$ kèm theo $\varphi$ tăng nhanh thường báo hiệu quá tải cảm kháng hoặc hỏng hóc cách điện.
Ngược lại, $S$ tăng trong khi $\varphi$ ổn định có thể chỉ ra quá tải thuần trở.
Bằng cách đưa các kênh phái sinh từ vật lý vào không gian đặc trưng,
các mô hình dự báo và phân loại phía sau sẽ có thêm các giả thiết tiên nghiệm
mà các đặc trưng thuần thống kê (trung bình trượt, năng lượng sóng) không thể cung cấp.

\subsection{Weather Integration}

Dữ liệu khí tượng từ Open-Meteo API~\cite{openmeteo2026} (tọa độ 34,94°N, 127,70°E)
được tích hợp để cung cấp ngữ cảnh môi trường cho tiêu thụ điện.
Do dữ liệu khí tượng gốc có tần suất 1 giờ trong khi dữ liệu điện là 15 phút,
chúng tôi thực hiện nội suy tuyến tính để resample về 15 phút:
\begin{equation}
w(t) = w(t_1) + \frac{w(t_2) - w(t_1)}{t_2 - t_1}(t - t_1), \quad t \in [t_1, t_2]
\label{eq:weather_interp}
\end{equation}
trong đó $w$ là biến khí tượng (nhiệt độ, độ ẩm, tốc độ gió).

Các đặc trưng khí tượng phái sinh được tính toán thêm:
\begin{itemize}
    \item \textbf{temp\_rolling\_24h:} Nhiệt độ trung bình 24 giờ (96 bước) để làm mượt biến động ngày-đêm.
    \item \textbf{temp\_diff\_24h:} Chênh lệch nhiệt so với cùng thờ điểm hôm trước, phản ánh xu hướng thờ tiết.
    \item \textbf{heat\_index:} Chỉ số nhiệt cảm nhận Rothfusz, proxy cho tải làm mát HVAC.
    \item \textbf{humidity\_x\_temp:} Tương tác nhiệt-ẩm, phản ánh điều kiện khí hậu khắc nghiệt.
\end{itemize}

Merge thực hiện bằng \texttt{left\_join} trên cột \texttt{date},
đảm bảo không làm thay đổi số dòng của dữ liệu điện và không sử dụng back-fill
nhằm tránh đưa thông tin thờ tiết tương lai vào quá khứ (zero data leakage).

\subsection{Proxy Anomaly Labeling Strategy}

Ba loại bất thường được định nghĩa bằng các ngưỡng dựa trên vật lý và được chuẩn hóa trong
\texttt{metadata/dataset/LABELING\_GUIDELINE.md}. Tài liệu guideline này đóng vai trò như artifact kiểm toán:
nêu nguồn cơ sở của từng rule, điều kiện gán nhãn, confidence score, ví dụ diễn giải,
và các lưu ý khi không có ground truth từ SCADA hoặc bảo trì.
\begin{itemize}
\item \textbf{Chạy không tải:} Tải nhẹ trong giờ nghỉ với tiêu thụ cao và PF $< 0{,}5$
\item \textbf{Rò rỉ:} Trung bình trượt 1 tuần vượt quá baseline 4 tuần đầu tiên 5\%
(theo ISO~50001 và Abraham et al.~\cite{abraham2021iso};
baseline 4 tuần giảm thiểu bias từ một tuần khởi động bất thường)
\item \textbf{Quá tải:} Tiêu thụ vượt quá phân vị 99{,}5 với công suất phản kháng cao và PF $< 0{,}7$
\end{itemize}

Mỗi nhãn được gán kèm confidence score trong $[0,1]$ phản ánh mức độ tin cậy của quyết định.
Score được tính bằng trung bình có trọng số của các chỉ báo nhị phân thành phần,
đảm bảo tính giải thích (explainability) và tránh ``black-box'' labeling.

\subsection{GAN-based Data Augmentation}

Mất cân bằng lớp là một thách thức nghiêm trọng trong phát hiện bất thường,
với tỷ lệ anomaly khoảng 6--7\%.
Chúng tôi sử dụng Mạng đối kháng Sinh (GAN) để tăng cường dữ liệu lớp thiểu số.
Kiến trúc được chọn là mạng fully connected tiêu chuẩn với hàm kích hoạt LeakyReLU
vì dữ liệu sau feature engineering là vector đặc trưng đa chiều,
không còn cấu trúc không gian-thời gian thô cần được bảo toàn bởi CNN hoặc RNN.


\textbf{Chiến lược huấn luyện:} Generator được huấn luyện \textbf{chỉ trên các mẫu bất thường} (minority class, 2.388 mẫu)
nhằm học phân phối đặc trưng của lớp thiểu số thay vì bị chi phối bởi lớp đa số.
Sau khi huấn luyện, Generator sinh ra 500 mẫu synthetic được nối vào tập dữ liệu gốc,
bổ sung thêm dữ liệu cho lớp anomaly nhằm hỗ trợ các bộ phân loại phía sau.
Kiến trúc FC-GAN được chọn làm baseline đơn giản;
tuy nhiên, theo Yoon et al.~\cite{yoon2019} và Li et al.~\cite{li2022ttsgan},
các kiến trúc chuyên biệt cho chuỗi thờ gian như TimeGAN hoặc TTS-GAN
sẽ bảo toàn động lực thờ gian tốt hơn cho dữ liệu công nghiệp.
Mục tiêu minimax của GAN được định nghĩa:
\begin{equation}
\begin{aligned}
\min_G \max_D \mathcal{L}(D,G)
&= \mathbb{E}_{x\sim p_{\text{data}}}\big[\ln D(x)\big] \\
&\quad + \mathbb{E}_{z\sim p_z}\big[\ln(1-D(G(z)))\big]
\end{aligned}
\label{eq:gan}
\end{equation}

Trong đó $G$ là generator nhận vector nhiễu $z \sim \mathcal{N}(0,\mathbf{I})$
và sinh mẫu tổng hợp $G(z)$; $D$ là discriminator phân biệt mẫu thật và giả.
Cả hai mạng được tối ưu bằng Adam với learning rate $2\times 10^{-4}$ và $\beta_1 = 0{,}5$.

\subsection{Pipeline Integration}

Pipeline tiền xử lý được triển khai dưới dạng luồng tuần tự gồm 7 bước (Alg.~\ref{alg:pipeline}).
Các kiểm tra khẳng định (assertions) được tích hợp sau mỗi bước
để đảm bảo tính tái lập và phát hiện lỗi tự động.

\begin{algorithm}[htbp]
\caption{DS108-ELECTRIMIGHT PIPELINE}
\label{alg:pipeline}
\begin{algorithmic}[1]
\REQUIRE Raw CSV $D_{\text{raw}}$, Weather CSV $W_{\text{raw}}$
\ENSURE Processed dataset $D_{\text{final}}$
\STATE $D_1 \leftarrow \text{load}(D_{\text{raw}})$ \COMMENT{Parse date với dayfirst=True}
\STATE $D_2, R_{\text{clean}} \leftarrow \text{clean}(D_1)$
\STATE \COMMENT{Remove duplicates, linear interpolation, PF scale, physical validation}
\STATE $D_3, R_{\text{weather}} \leftarrow \text{integrate\_weather}(D_2, W_{\text{raw}})$
\STATE \COMMENT{Resample 1h→15min, engineer derived features, left join}
\STATE $D_4 \leftarrow \text{build\_time\_features}(D_3)$
\STATE \COMMENT{Lag $\{1,2,4,96\}$, rolling $\{24,48,96\}$, cyclical NSM}
\STATE $D_5 \leftarrow \text{rolling\_wavelet\_features}(D_4, \text{db4}, L=3, w=64)$
\STATE $D_6 \leftarrow \text{build\_physical\_features}(D_5)$
\STATE \COMMENT{Compute $S = \sqrt{P^2+Q^2}$ and $\varphi = \arccos(\text{PF})$}
\STATE $D_7 \leftarrow \text{label\_all\_anomalies}(D_6)$
\STATE \COMMENT{Apply idling/leakage/overload rules and confidence scores}
\STATE $\text{run\_pipeline\_assertions}(D_7, \text{stage}=\text{post\_labeling})$
\IF{use\_GAN}
    \STATE $G, D_{\text{net}} \leftarrow \text{train\_GAN}(D_7^{(\text{anomaly})})$
    \STATE $D_{\text{syn}} \leftarrow \text{generate}(G, n_{\text{syn}})$
    \STATE $D_{\text{final}} \leftarrow D_7 \cup D_{\text{syn}}$
    \STATE $\text{assert\_correlation\_preserved}(D_7, D_{\text{syn}})$
\ELSE
    \STATE $D_{\text{final}} \leftarrow D_7$
\ENDIF
\STATE \textbf{return} $D_{\text{final}}$
\end{algorithmic}
\end{algorithm}

% =============================================================================
